\scanpage{89}
% archive foliation: 88

$2$ forment les sections d'un torseur sous $(\unsure{\mathfrak{S}|\mathcal G}_3)_S$. Donc les bi-épinglages forment les sections d'un torseur sous le groupe \deleted{\unsure{étale}}
\[
(45)\qquad(\unsure{\boldsymbol{\mu}|\mu}_4\times\unsure{\mathfrak{S}|\mathcal G}_3)_S
\]
donc on trouve une action\inserted{à g.} du groupe absolu $(\unsure{\boldsymbol{\mu}|\mu}_4\times\unsure{\mathfrak{S}|\mathcal G}_3)_{S_0}$ sur \deleted{\unsure{X}} le schéma modulaire (37), $U_{0,3}\times_{S_0}\unsure{\boldsymbol{\mu}|\mu}_4^*$, qu'il faudrait expliciter, en même temps que l'action sur
\[
(46)\qquad\mathbb{T}_0=\underline{O}_{U_{0,3}\times\unsure{\boldsymbol{\mu}|\mu}_4^*}(1)
\]
à l'aide de l'action ``tautologique'' de $\unsure{\mathfrak{S}|\mathcal G}_3$ sur \deleted{\unsure{$\mathbb{P}$}} $U_{03}=\Sigma_{\{0,1,\infty\}}\setminus\{0,1,\infty\}_{S_0}$ via l'action de $\unsure{\mathfrak{S}|\mathcal G}_3$ sur $V_{\deleted{\unsure{J}}\{0,1,\infty\}}(\underline{O}_{S_0})$.
\[
(47)\quad\left\{\begin{array}{l}
\unsure{\mathfrak{S}|\mathcal G}_3\text{ opère sur }U_{03}\times\unsure{\boldsymbol{\mu}|\mu}_4^*\text{ via }\left\{\begin{array}{l}\text{action tautol. sur }U_{03}\\\text{le caractère }\operatorname{sg}(\sigma)\text{ sur }\unsure{\boldsymbol{\mu}|\mu}_4^*\end{array}\right.\\[7pt]
\unsure{\boldsymbol{\mu}|\mu}_4\text{ opère sur }U_{03}\times\unsure{\boldsymbol{\mu}|\mu}_4^*\text{ via }\left\{\begin{array}{l}\text{action triviale sur }U_{03}\\\text{le car. can. }\unsure{\boldsymbol{\mu}|\mu}_4\to\unsure{\boldsymbol{\mu}|\mu}_2=\{\pm1\}_{S_0}\\\text{sur }\unsure{\boldsymbol{\mu}|\mu}_4^*\end{array}\right.
\end{array}\right.
\]
donc $\unsure{\mathfrak{S}|\mathcal G}_3\times\unsure{\boldsymbol{\mu}|\mu}_4$ opère sur le facteur $U_{03}$ via l'
